\section{Data-Centric Validation and the Limits of Synthetic Benchmarks}
The study treats synthetic narratives as a controlled experimental instrument rather than as a statistically representative sample of barangay residents. This position requires both a defense and a limitation. Long et al. \cite{long} organize synthetic-data research around generation, curation, and evaluation rather than treating large-scale generation as sufficient evidence of quality. Chim et al. \cite{chim} emphasize examining the differences between source text and generated text, including style and downstream usefulness. In this study, repeatable generation makes it possible to vary language, issue, and ambiguity by design, but repeatability does not automatically ensure semantic accuracy or transfer to unfamiliar writing.

The possibility of \emph{annotation artifacts} is particularly important. Gururangan et al. \cite{gururangan} demonstrated that superficial regularities introduced by dataset construction can support shortcuts in a natural-language task. Shaib et al. \cite{shaib} study recurring syntactic templates in generated text. These findings motivate three distinct checks in the present research: (a) strict non-overlap of recorded scenario-family identifiers between partitions; (b) measurement of text-length and lexical proxies for label identity; and (c) linguistic tests beyond the generator-derived validation distribution. Group separation prevents the same registered scenario family from occurring in both sides of a split, but it does not eliminate shared words or class-specific narrative frames produced by the same generator.

\section{Why Cross-Validation and a Separate Validation Set Both Matter}
A single train--test split can yield a score that depends on which scenarios happened to be selected for evaluation. Cross-validation repeats training across multiple held-out subsets and therefore exposes fold-to-fold variation. However, ordinary K-fold splitting allows closely related variations of one generating scenario to appear on both sides. The study instead employs \texttt{StratifiedGroupKFold}, which attempts to maintain class proportions while ensuring that each group appears in only one held-out fold \cite{sklearngroup}. Importantly, the TF-IDF vectorizer is inside the scikit-learn pipeline; its document frequencies and vocabulary are learned from the training portion of each fold, never from the held-out fold \cite{sklearnpitfalls}.

The three-fold grid search is used for \emph{selection} of the SVM parameters. A subsequent five-fold group-based analysis is reported as a stability diagnostic for the selected SVM and the fixed Naive Bayes baseline. The fixed validation partition is then assessed separately. These measures address different questions: the grid compares alternatives; the second set of folds inspects repeated partition behavior; and validation measures performance on the designated sample that was not used to fit model weights. The reserved final test partition is \emph{not} used for the present study tuning or the reported validation experiments.

\section{Behavioral Testing Beyond Aggregate Accuracy}
Ribeiro et al. \cite{ribeiro} argue that aggregate accuracy can miss systematic linguistic weaknesses and present the CheckList approach for behavioral tests such as changes that should preserve or alter a prediction. This suggests that the final application needs controlled cases involving paraphrases, negation, incomplete references, code-switching, and urgent but informal reports. Our the present study validation examines assigned-language subgroups and simple surface-feature baselines; it does not claim that a complete, independently adjudicated behavioral benchmark has been established. The distinction is necessary to avoid presenting promising model experiments as production validation.

\section{Documentation and Independent Annotation}
Gebru et al. \cite{datasheets} propose datasheets that describe the origin, composition, intended uses, and limitations of datasets. We adapt this principle by retaining metadata for each generated example: source identifiers, classification rationale, intended language, scenario-family membership, generation method, and partition. Artstein and Poesio \cite{artstein} explain why annotation agreement must be interpreted in context rather than treated as interchangeable with classifier accuracy. A future blind review with independent raters is therefore a separate requirement, and the present records are explicitly flagged as not yet human-reviewed. These caveats are reflected in the scope and interpretation of Chapter 4.
